\documentclass[11pt,a4paper]{article}
\usepackage[utf8]{inputenc}
\usepackage[T1]{fontenc}
\usepackage[ngerman]{babel}
\usepackage{lmodern}
\usepackage[a4paper,margin=2.35cm]{geometry}
\usepackage{booktabs,tabularx,longtable,array,enumitem,graphicx,float,parskip}
\usepackage[hidelinks]{hyperref}
\newcolumntype{Y}{>{\raggedright\arraybackslash}X}
\setlist[itemize]{leftmargin=1.4em,itemsep=2pt,topsep=3pt}
\setlist[enumerate]{leftmargin=1.6em,itemsep=2pt,topsep=3pt}
\setlength{\emergencystretch}{3em}
\newcommand{\appscreenshot}[3]{\IfFileExists{#1}{%
\begin{figure}[H]\centering
\includegraphics[width=.94\linewidth,height=.56\textheight,keepaspectratio]{#1}
\caption{#2}\label{#3}
\end{figure}}{}}

\begin{document}
\begin{titlepage}\centering
\vspace*{2cm}
{\Large Technische Hochschule Mittelhessen\par}
{\large Fachbereich MNI\par}
{\large Praktikum Wirtschaftsinformatik, Sommersemester 2026\par}
\vspace{2.5cm}\rule{\linewidth}{.8pt}\par\vspace{.5cm}
{\Huge\bfseries Aufgabendatenbank\par}
{\LARGE Projektdokumentation\par}
\vspace{.5cm}\rule{\linewidth}{.8pt}\par\vspace{2cm}
\begin{tabular}{rl}
\textbf{Betreuung:} & Prof. Dr. Markus Siepermann; Johannes Kunz\\
\textbf{Projektteam:} & [Name 1] (Projektleitung)\\
& [Name 2]\\ & [Name 3]\\
\textbf{Gruppe:} & [Gruppennummer]\\
\textbf{Abgabe:} & 14. Juli 2026
\end{tabular}
\vfill{\small Repository: \texttt{PWI-Aufgabendatenbank}}
\end{titlepage}

\tableofcontents\newpage

\section{Kurze Zusammenfassung}
Ausgangspunkt war ein hochschulweit abgestimmtes relationales Schema zur
Beschreibung und Organisation von Übungsaufgaben, für das noch keine
nutzbare Anwendung existierte. Das Projekt überführt dieses abstrakte Modell
in ein Autorenwerkzeug aus Vue-Frontend, Spring-Boot-Backend und
PostgreSQL-Datenbank.

Die Anwendung verwaltet Aufgaben mit typisierten Inhaltsbausteinen,
Metadaten, Darstellungstemplates und Validatoren. Aufgaben können als
ungeordnete Sammlungen, geordnete Sequenzen oder Varianten derselben
Ausgangsaufgabe organisiert werden. Hierarchische Tags und kombinierte Filter
unterstützen das Wiederfinden. Das Ergebnis ist kein reines CRUD-System:
Beziehungen wie \texttt{purpose}, \texttt{position} und
\texttt{root\_item\_id} wurden in konkrete, konsistente Bedienabläufe
übersetzt. Automatisierte Tests, Docker und eine GitHub-Actions-Pipeline
ergänzen die Anwendung.

\section{Aufgabenstellung und Zielsetzung}
Die Aufgabe bestand nicht darin, ein eigenes vereinfachtes Datenmodell zu
entwerfen. Das vorgegebene Schema mit Entitäten wie \texttt{item},
\texttt{item\_content}, \texttt{item\_collection}, \texttt{tag},
\texttt{validator} und \texttt{modifier} sollte fachlich verstanden,
technisch abgebildet und für reale Nutzende bedienbar gemacht werden.

Zielgruppe sind in erster Linie Lehrende und Autoren von Übungsaufgaben. Sie
sollen Aufgaben erstellen, mit strukturierten Inhalten und Metadaten
ausstatten, kombinieren und wiederfinden können. Die Anwendung ist weder ein
Lernmanagementsystem noch eine Laufzeitumgebung zur Abgabe und automatischen
Korrektur. Validatoren und Modifier werden als Regeln modelliert; ihre
Ausführung gehört zu einem späteren externen Dienst.

Die Projektziele waren:
\begin{itemize}
\item vollständige Abbildung der zentralen Schemaelemente ohne neue Tabellen,
\item verständliche Arbeitsabläufe für Aufgaben, Inhalte und Beziehungen,
\item konsistente REST-Schnittstellen zwischen Frontend und Backend,
\item persistente Speicherung in PostgreSQL,
\item automatisierte Prüfung kritischer Geschäftsregeln,
\item technische Vorbereitung auf die übergeordnete Plattform.
\end{itemize}

\section{Ausgangssituation und Recherche}
Zu Projektbeginn lagen das Schema, ein Lastenheft und ein Vue-basiertes
Frontend-Gerüst der übergeordneten EWiLL-Plattform vor. Viele Beziehungen
waren technisch benannt, ihre Bedeutung für die Oberfläche war jedoch nicht
vollständig festgelegt. Das Team analysierte deshalb zuerst Primär- und
Fremdschlüssel, Kardinalitäten und Beziehungsattribute und klärte offene
Punkte in Betreuerterminen.

Besonders relevant waren folgende Ergebnisse:
\begin{itemize}
\item Ein \texttt{item} trägt Aufgabenmetadaten; der eigentliche Text oder
eine Datei liegt in einem oder mehreren \texttt{item\_content}-Datensätzen.
\item \texttt{purpose} gehört zur Verbindung
\texttt{item\_contents}, weil es die Rolle eines Inhalts in einer konkreten
Aufgabe beschreibt.
\item Collections modellieren die vertikale Gruppierung beziehungsweise
Sequenz; ihre Reihenfolge liegt in
\texttt{item\_collection\_sub\_item.position}.
\item \texttt{root\_item\_id} ist keine Baum-Elternreferenz, sondern verbindet
horizontale Varianten derselben Ausgangsaufgabe.
\item Validatoren und Modifier sind wiederverwendbare Beschreibungen. Die
Aufgabendatenbank speichert und verknüpft sie, führt sie aber nicht aus.
\end{itemize}

\subsection{Technologieauswahl}
Die Auswahl beruhte auf fachlicher Eignung, vorhandener Teamerfahrung und den
Integrationsbedingungen der übergeordneten Plattform. Da das bereitgestellte
Frontend-Gerüst bereits Vue verwendete und das Team damit vertraut war, wurde
es erweitert statt durch ein zweites Framework ersetzt. Auch mit Spring Boot
und PostgreSQL bestand bereits Erfahrung. So konnte der Schwerpunkt auf
Datenmodell und Bedienlogik liegen.

\begin{longtable}{|p{.18\linewidth}|p{.28\linewidth}|p{.45\linewidth}|}
\caption{Verwendete Technologien; Versionen gemäß finaler Projektkonfiguration}
\label{tab:tech}\\
\hline
\textbf{Bereich}&\textbf{Technologie und Version}&\textbf{Einsatz und Begründung}\\
\hline
\endfirsthead
\hline
\textbf{Bereich}&\textbf{Technologie und Version}&\textbf{Einsatz und Begründung}\\
\hline
\endhead
Frontend&Vue 3.5.34, Vue Router 5.0.6, TypeScript 6.0.3, Axios 1.16.0&
Erweiterung des bereitgestellten Vue-Frontends; komponentenbasierte,
typsichere Oberfläche und REST-Kommunikation. Das Team hatte bereits
Vue-Erfahrung.\\
UI und Interaktion&Vuetify 4.0.7, Pinia 3.0.4, vuedraggable 4.1.0,
CodeMirror 6&Einheitliche Bedienelemente, zentraler Zustand, Drag-and-Drop
sowie XML-Bearbeitung der Darstellungstemplates.\\
Backend&Java 21, Spring Boot 3.2.5, Spring Data JPA&REST-Schnittstellen,
Geschäftslogik, Transaktionen und Abbildung des relationalen Schemas. Spring
Boot war dem Team bekannt und für die Schichtenarchitektur geeignet.\\
Datenbank&PostgreSQL 16 Alpine&Native Unterstützung der im Schema verwendeten
UUIDs, JSONB-Inhalte, GIN-Indizes und BYTEA-Dateien.\\
Tests und Build&Vitest 4.1.6, JUnit 5, Mockito, MockMvc, Vite 8.0.12,
Maven Wrapper 3.9.6&Automatisierte Prüfung von Frontend, Services und
Controllern sowie reproduzierbare Produktionsbuilds.\\
Betrieb und Plattform&Docker Compose, GitHub Actions, GHCR, Kubernetes/ArgoCD&
Reproduzierbarer lokaler Start und automatischer Image-Build. Das
Cluster-Deployment wurde vorbereitet, blieb aber von externer Infrastruktur
abhängig.\\
\hline
\end{longtable}

Wichtig ist die genaue Bezeichnung: Eingesetzt wird \textbf{Vue 3.5.34}.
Die Versionsnummer 5.0.6 gehört zu \textbf{Vue Router}, nicht zu Vue selbst.

\section{Anforderungen und Rahmenbedingungen}
\subsection{Funktionale Anforderungen}
\begin{enumerate}[label=F\arabic*]
\item Aufgaben mit Autor, Lizenz, Typ und optionalem Template verwalten.
\item Inhaltsbausteine mit Zweck, Typ, Autor und Lizenz verwalten.
\item Text, strukturierte Daten, Bilder und PDF-Dateien speichern und anzeigen.
\item Geordnete Sequenzen und ungeordnete Sammlungen bilden.
\item Varianten einer Ausgangsaufgabe verwalten.
\item Validatoren definieren und Aufgaben zuordnen.
\item Hierarchische Tags anlegen, zuordnen und filtern.
\item Nach Text, Autor, Typ und Tags suchen.
\item Darstellungstemplates bearbeiten und als Vorschau anwenden.
\end{enumerate}

\subsection{Nicht-funktionale Anforderungen}
\begin{itemize}
\item referentielle Integrität und transaktionale Geschäftsoperationen,
\item typisierte Verträge zwischen Frontend und Backend,
\item verständliche Bedienung trotz abstraktem Datenmodell,
\item automatisierte Tests für kritische Erfolgs- und Fehlerfälle,
\item reproduzierbarer Start in Containern,
\item keine strukturelle Erweiterung des vorgegebenen Schemas.
\end{itemize}

\section{Konzept und Entwurf}
\subsection{Systemarchitektur}
Die Anwendung folgt einer Schichtenarchitektur:
\begin{verbatim}
Vue-Komponente -> Pinia Store -> ApiAdapter -> Axios -> REST-Controller
                                                   -> Service
                                                   -> Repository
                                                   -> PostgreSQL
\end{verbatim}
Controller definieren HTTP-Vertrag und Statuscodes. Services validieren
Referenzen, kapseln Geschäftsregeln und Transaktionsgrenzen. Spring-Data-
Repositories übernehmen den Datenzugriff. DTOs bilden den API-Vertrag, sodass
JPA-Entities mit Lazy Loading und Persistenzdetails nicht direkt serialisiert
werden. Im Frontend übersetzt der Store DTOs in UI-Modelle. Die Adapter-
Schnittstelle entkoppelt ihn von Axios und erlaubt einen lokalen Dummy-Betrieb.

\subsection{Beziehungsmodell}
\begin{table}[H]\centering
\begin{tabularx}{\linewidth}{|l|Y|Y|}\hline
&\textbf{Vertikal}&\textbf{Horizontal}\\\hline
Bedeutung&Gruppierung verschiedener Aufgaben oder Lernsequenz&
Varianten derselben Ausgangsaufgabe\\
Umsetzung&\texttt{item\_collection},
\texttt{item\_collection\_sub\_item}, optional \texttt{position}&
\texttt{item.root\_item\_id}, ergänzend Validatoren und Modifier\\
Reihenfolge&optional; ausschließlich in der Join-Tabelle&keine\\
\hline
\end{tabularx}
\caption{Fachliche Trennung der Beziehungsarten}
\end{table}

\subsection{Konsistenzentscheidungen}
Beziehungen mit eigenen Attributen wurden als eigenständige JPA-Entities mit
zusammengesetztem Schlüssel modelliert. Das betrifft
\texttt{item\_contents} mit \texttt{purpose} und
\path{item_collection_sub_item} mit \texttt{position}. Ein einfaches
\texttt{@ManyToMany} hätte diese fachlichen Attribute verloren.

Das Backend ist die Autorität für Reihenfolgen. Beim Aktivieren einer
Sortierung erzeugt es fortlaufende Positionen, beim Deaktivieren setzt es alle
Positionen auf \texttt{NULL}, und nach einem Verschieben normalisiert es die
Geschwister. Die Oberfläche arbeitet bei mehrstufigen Vorgängen optimistisch;
nach einem Teilfehler lädt sie den betroffenen Aggregatbereich erneut. Dadurch
bleiben Datenbank und Oberfläche konsistent.

\section{Umsetzung}
\subsection{Abläufe aus Sicht der Nutzenden}
\textbf{Initiales Laden.} Beim Öffnen lädt das Frontend Referenzdaten wie
Autoren, Lizenzen und Typen sowie die Wurzelstruktur. Für ausgewählte Items
werden Inhalte und Beziehungen nachgeladen. Die flachen API-Daten werden im
Pinia-Store in die Baum- und Editoransicht übersetzt.

\textbf{Aufgabe und Inhalt erstellen.} Der Dialog prüft die Eingaben und sendet
zuerst \path{POST /api/items}. Die zurückgegebene UUID wird anschließend für
\path{POST /api/contents/by-item/itemId} verwendet. Der Backend-Service
prüft Autor, Lizenz und Typ, speichert den Content und erzeugt die Verbindung
mit ihrem \texttt{purpose}. Bei einem partiellen Fehler synchronisiert das
Frontend erneut mit dem Server.

\textbf{Collection bearbeiten.} Ein sichtbarer Collection-Knoten besteht aus
einem Träger-Item und einem eigenen Collection-Datensatz. Beim Drag-and-Drop
sendet das Frontend Quell- und Zielkontext. Der Service ändert die
Mitgliedschaft und berechnet bei geordneten Collections die Positionen neu.
Das Variantenfeld \texttt{rootItemId} bleibt dabei unverändert. Das Löschen
einer Collection entfernt nicht automatisch ihre bisher enthaltenen Aufgaben.

\textbf{Variante erstellen.} Eine neue Variante ist ein eigenständiges Item,
dessen \texttt{root\_item\_id} auf die Ausgangsaufgabe verweist. Sie besitzt
eigene Inhalte und Metadaten. Das Variantenpanel lädt und zeigt alle Items
derselben Aufgabenfamilie unabhängig von ihrer Collection-Zugehörigkeit.

\textbf{Tags, Validatoren und Templates.} Tag- und Validator-Zuordnungen
werden über eigene REST-Endpunkte gespeichert. Das Frontend baut aus der
flachen Tag-Liste rekursiv einen Baum. Templates werden als XML-Text
gespeichert; das Frontend liest die Reihenfolge der
\texttt{purpose}-Platzhalter und ordnet damit die Inhaltsvorschau.

\appscreenshot{docs/screenshots/01-aufgabe-erstellen.png}
{Erstellung einer Aufgabe mit zentralen Metadaten.}{fig:erstellung}
\appscreenshot{docs/screenshots/02-inhalte-und-vorschau.png}
{Inhaltsbausteine und templatebasierte Vorschau.}{fig:inhalte}
\appscreenshot{docs/screenshots/03-kollektionen.png}
{Geordnete Sequenz und ungeordnete Sammlung.}{fig:collections}
\appscreenshot{docs/screenshots/04-varianten-validatoren.png}
{Varianten und zugeordnete Validatoren.}{fig:varianten}
\appscreenshot{docs/screenshots/05-tags-und-suche.png}
{Tag-Baum sowie kombinierte Suche und Filter.}{fig:tags}

\subsection{Umsetzung aller Entitäten}
Tabelle~\ref{tab:entities} trennt bewusst Datenbank, Backend und Frontend.
Der Status teilweise bedeutet, dass das Schemaelement berücksichtigt ist,
aber kein vollständiger eigener Bedienablauf vorliegt.

\small
\begin{longtable}{|p{.13\linewidth}|p{.24\linewidth}|p{.28\linewidth}|p{.24\linewidth}|}
\caption{Umsetzung der Entitäten auf allen Ebenen}\label{tab:entities}\\
\hline
\textbf{Entität}&\textbf{Datenbank und Bedeutung}&\textbf{Backend}&
\textbf{Frontend und Status}\\\hline
\endfirsthead
\hline
\textbf{Entität}&\textbf{Datenbank und Bedeutung}&\textbf{Backend}&
\textbf{Frontend und Status}\\\hline
\endhead
\texttt{author}&UUID, Bezeichnung und E-Mail; verantwortlich für Item oder
Content, aber kein Benutzerkonto.&JPA-Entity, Repository und Referenz-
Controller; Schreiboperationen prüfen die ID.&Auswahl und Neuanlage in den
Metadaten. Für die erste Iteration vollständig umgesetzt.\\
\texttt{license}&Wiederverwendbare Lizenz; Item und Content dürfen eigene
Lizenzen besitzen. Fremdschlüssel verhindern das Löschen verwendeter
Einträge.&Entity, Repository, Lese- und Erstellendpunkte,
Referenzvalidierung.&Auswahl und Neuanlage für Item beziehungsweise Content.
Umgesetzt.\\
\texttt{item\_type}&Klassifiziert eine Aufgabe. Verbindung zu erlaubten
Content-Typen über \texttt{item\_content\_types}.&Entity, Repository und
Referenzendpunkte; JPA-Beziehung zu Content-Typen.&Auswahl, Neuanlage und
Suchfilter umgesetzt; automatische Einschränkung kompatibler Content-Typen
noch nicht vollständig in der UI.\\
\path{item_content_type}&Beschreibt Format beziehungsweise fachliche Art
eines Inhaltsblocks.&Entity, Repository, Referenzendpunkte und Prüfung im
Content-Service.&Auswahl und Neuanlage im Content-Editor. Umgesetzt.\\
{\tiny\path{item_representation_template}}&Optional vom Item referenzierter
XML-Template-Text; trennt Inhalte von ihrer Anordnung.&Entity, Repository,
Service und CRUD-Controller.&CodeMirror-Editor, sortierbare Purpose-
Bausteine, Speichern und Vorschau. Umgesetzt als projektspezifische erste
Template-Sprache.\\
\texttt{tag}&Hierarchisches Schlagwort mit Self-Reference
\texttt{parent\_tag\_id}; beim Löschen bleibt ein Kind als Wurzel erhalten.&
Entity, Repository und Controller; flache DTO-Liste, Anlage und Löschen.&
Rekursiver Baum, Pfade, Zuordnung, Entfernung und Filter. Item-Tags umgesetzt;
eigener Content-Tag-Workflow offen.\\
\texttt{validator}&Wiederverwendbare Prüf- oder Restriktionsbeschreibung.&
Entity, Repository, Service, CRUD und Endpunkte für die Item-Zuordnung.&
Editor sowie Zuordnen und Entfernen an Aufgaben. Speicherung umgesetzt,
Ausführung bewusst extern.\\
\texttt{modifier}&Regel zur Transformation oder Erzeugung einer Variante.&
Entity und Repository vorhanden; Item-Mapping und IDs im DTO/Service
berücksichtigt.&Kein vollständiger Modifier-Editor und keine Ausführung.
Strukturell vorbereitet, teilweise umgesetzt.\\
\texttt{item}&Zentrale Aufgabenentität mit Autor, Lizenz, Typ, optionalem
Template und \texttt{root\_item\_id} für Varianten.&Entity, Repository,
CRUD-Service/-Controller, Suche, Varianten-, Tag- und Validatoroperationen.&
Baumknoten, Editor, Metadaten, Inhalte, Tags, Validatoren und Varianten.
Kernfunktion umgesetzt.\\
\texttt{item\_\allowbreak content}&Eigenständiger Content mit Autor, Lizenz, Typ, JSONB
oder BYTEA sowie Zeitstempeln.&Entity, Repository, Service und Controller;
CRUD, Multipart-Upload und Blob-Download.&Content-Liste und Editor; Text,
Bildvorschau und Datei-Download. Kernworkflow umgesetzt; Wiederverwendung
desselben Contents in mehreren Items nicht als eigener UI-Workflow.\\
\texttt{item\_\allowbreak collection}&Eigene UUID, Träger-Item und
\texttt{order}-Kennzeichen.&Entity, Repository, transaktionaler Service und
Controller; Konvertierung ist idempotent.&Darstellung als Ordner mit
\texttt{collectionId}, Kindern und Sortierstatus. Umgesetzt.\\
\texttt{item\_\allowbreak collection\_\allowbreak sub\_\allowbreak item}&Mitgliedschaft zwischen Collection und
Item; \texttt{position} ist nur bei Sequenzen gesetzt.&Eigene Entity mit
Composite Key; Hinzufügen, Entfernen, Verschieben und Normalisieren.&
Drag-and-Drop mit vuedraggable; Reihenfolge kommt abschließend vom Backend.
Umgesetzt.\\
\hline
\end{longtable}
\normalsize

\subsection{Umsetzung der Join-Tabellen}
\begin{longtable}{|p{.24\linewidth}|p{.68\linewidth}|}
\caption{Bedeutung und Status der Join-Tabellen}\label{tab:joins}\\
\hline
\textbf{Tabelle}&\textbf{Umsetzung}\\\hline
\endfirsthead
\hline
\textbf{Tabelle}&\textbf{Umsetzung}\\\hline
\endhead
\texttt{item\_contents}&Verbindet Item und Content und trägt
\texttt{purpose}. Deshalb eigene JPA-Entity mit Composite Key. Backend und
Content-Editor verwenden die Beziehung vollständig.\\
\texttt{item\_\allowbreak collection\_\allowbreak sub\_\allowbreak item}&Verbindet Collection und Item und trägt
\texttt{position}. Eigene JPA-Entity; Service und Drag-and-Drop verwenden sie
vollständig.\\
\texttt{item\_tags}&Many-to-many-Zuordnung ohne Zusatzattribute. JPA-Set,
Zuordnungsendpunkte, Chips, Baumwahl und Filter sind umgesetzt.\\
\texttt{item\_content\_tags}&In Entity, DTO und Service berücksichtigt.
Ein gleichwertiger eigener Content-Tag-Editor wurde nicht fertiggestellt.\\
\texttt{item\_validator}&Many-to-many-Verbindung; Backend-Endpunkte und
Validatorpanel sind für Speicherung und Zuordnung umgesetzt.\\
\texttt{item\_modifier}&Item-Zuordnung im Datenmodell und Backend-Mapping
vorbereitet; vollständiges CRUD und UI fehlen.\\
\texttt{item\_content\_types}&Kompatibilitätsbeziehung zwischen Aufgaben- und
Content-Typ. In SQL und JPA modelliert; automatische UI-Filterung ist offen.\\
\hline
\end{longtable}

\subsection{Suche und Bedienlogik}
Der Item-Endpunkt akzeptiert optionale Kriterien. Das Backend kombiniert sie
über JPA Specifications und bezieht JSONB-Inhalte in die Textsuche ein. Das
Frontend kombiniert Suchtext, Autor, Aufgabentyp und Tag-Auswahl und kann die
Filter gemeinsam zurücksetzen. Hierarchische Tags werden von der API flach
übertragen und erst im Frontend zum Baum aufgebaut. Dadurch bleibt der
REST-Vertrag einfach und die Darstellung flexibel.

\subsection{Qualitätssicherung}
Zum finalen Stand liefen 107 Frontend- und 92 Backend-Tests erfolgreich.
Frontend-Tests prüfen unter anderem Store-Aktionen, DTO-Transformation,
Baumvalidierung, Collections, Varianten, Suche und Resynchronisation nach
partiellen API-Fehlern. Backend-Tests prüfen Item-, Content-, Collection-,
Validator- und Template-Services. MockMvc-Tests kontrollieren zentrale
Endpunkte und HTTP-Statuscodes. Zusätzlich waren TypeScript-Typcheck und
Vite-Produktionsbuild erfolgreich.

Der frühere leere Spring-Kontexttest wurde entfernt, weil er ohne fachliche
Assertion lediglich eine lokal laufende PostgreSQL-Instanz verlangte. Die
fachlichen Service- und Controller-Tests bleiben datenbankunabhängig. Ein
vollständiger Integrationstest gegen reales PostgreSQL, beispielsweise mit
Testcontainers, ist eine sinnvolle nächste Ausbaustufe.

\subsection{Build und Bereitstellung}
Docker Compose startet PostgreSQL, Backend und Frontend gemeinsam. GitHub
Actions baut nach Push auf \texttt{main} je ein Backend- und Frontend-Image
und veröffentlicht beide mit \texttt{latest}- und Commit-SHA-Tag in GHCR. Das
SHA-Tag ermöglicht einen reproduzierbaren Deploymentstand.

Eine Kubernetes-\texttt{values.yaml} wurde als Integrationsvorbereitung
erstellt. Das geplante produktive Deployment über ArgoCD konnte nicht
abgeschlossen werden: Die dafür zuständige Plattformgruppe führte ihr Projekt
nicht weiter und konnte keine nutzbare Cluster-/ArgoCD-Integration
bereitstellen. Dies ist eine externe Abhängigkeit, keine lokal lösbare
Anwendungsfunktion.

\section{Projektorganisation}
\subsection{Arbeitsweise mit GitHub}
GitHub war nicht nur Quellcodeablage, sondern das zentrale
Koordinationswerkzeug. Anforderungen und Fehler wurden als Arbeitspakete in
einem GitHub Project erfasst und abhängig vom Arbeitsstand zwischen
\emph{In progress}, \emph{In review} und \emph{Closed/Done} verschoben.
Dadurch waren Zuständigkeit und Fortschritt für das gesamte Team sichtbar.

Die Umsetzung erfolgte auf Feature- und Fix-Branches. Pull-Requests machten
Änderungen nachvollziehbar und ermöglichten Review vor dem Merge. Für
übergreifende Änderungen an Store, API-Adapter und DTOs wurden Tests und
Typprüfung vor der Integration ausgeführt. Merge-Konflikte traten insbesondere
bei parallel bearbeiteten zentralen Dateien auf; kleinere Arbeitspakete und
häufigere Integration reduzierten dieses Risiko.

\appscreenshot{docs/screenshots/06-github-project.png}
{GitHub Project mit Arbeitsständen und zugeordneten Arbeitspaketen.}
{fig:github-project}
\appscreenshot{docs/screenshots/07-pull-request.png}
{Beispiel eines Pull-Requests mit Review und erfolgreicher Prüfung.}
{fig:pull-request}

\subsection{Kommunikation und Zusammenarbeit}
Das Team traf sich abhängig von Fortschritt und Klärungsbedarf ein- bis
zweimal pro Woche. In diesen Terminen wurden offene fachliche Fragen,
Schnittstellen zwischen Frontend und Backend, Blockaden und die nächsten
Arbeitspakete abgestimmt. Zusätzliche kurze Absprachen erfolgten bei
Integrationsproblemen.

Einige komplexe Stellen wurden in gemeinsamen Pair-Programming-Sitzungen
bearbeitet. Dazu gehörten insbesondere schichtübergreifende Abläufe, bei denen
Frontend-Modell, REST-DTO und Backend-Regel gleichzeitig verstanden werden
mussten. Pair Programming diente hier sowohl der schnelleren Fehleranalyse als
auch dem Wissenstransfer innerhalb des Teams.

\subsection{Projektphasen und Meilensteine}
\begin{table}[H]\centering
\begin{tabularx}{\linewidth}{|l|l|Y|}\hline
\textbf{Phase}&\textbf{Zeitraum}&\textbf{Ergebnis}\\\hline
Analyse&Ende April--Mai&Schemaanalyse, Abgrenzung, Technologieentscheidung\\
Grundaufbau&Mai&Datenbank, Backend- und Frontend-Grundstruktur, Docker\\
Kernfunktionen&Juni&Items, Contents, Collections und REST-Integration\\
Erweiterungen&Juni--Juli&Varianten, Validatoren, Templates, Tags und Suche\\
Stabilisierung&Juli&Tests, Fehlerbehandlung, CI/CD und Dokumentation\\
Abgabe&14. Juli&Bereinigter Stand, Dokumentation und Präsentation\\
\hline
\end{tabularx}
\caption{Projektphasen}
\end{table}

\subsection{Aufgabenverteilung}
\begin{table}[H]\centering
\begin{tabularx}{\linewidth}{|l|Y|}\hline
\textbf{Person}&\textbf{Nachweisbarer Schwerpunkt}\\\hline
{[Name 1]}&[Projektleitung, Koordination und konkrete Beiträge]\\
{[Name 2]}&[Konkrete Frontend-, Backend- oder Datenbankbeiträge]\\
{[Name 3]}&[Konkrete Frontend-, Test- oder Deploymentbeiträge]\\
\hline
\end{tabularx}
\caption{Aufgabenverteilung; vor Abgabe anhand des GitHub Projects und der
Pull-Requests zu konkretisieren}
\end{table}

\subsection{Abweichungen vom ursprünglichen Plan}
Das Team wurde im Projektverlauf von fünf auf drei aktive Personen reduziert.
Der Umfang musste deshalb priorisiert und Aufgaben mussten
schichtübergreifend neu verteilt werden.

Mehrere im Lastenheft vorgesehene Punkte wurden bewusst nicht oder nur
teilweise umgesetzt:
\begin{itemize}
\item \textbf{KI-Unterstützung:} Die im Lastenheft erwogene KI-Funktion wurde
zugunsten eines konsistenten Kernsystems aus Datenmodell, Editor und
Organisation zurückgestellt. Es fehlte außerdem eine abschließend definierte
fachliche Aufgabe, anhand derer Nutzen und Qualität hätten geprüft werden
können.
\item \textbf{Modifier:} Datenbank, Entity, Repository und Item-Zuordnung sind
strukturell vorbereitet. Ein vollständiger Editor sowie eine
Ausführungslogik wurden nicht umgesetzt, weil die Laufzeittransformation zu
einem externen Dienst gehört und Kernfunktionen priorisiert wurden.
\item \textbf{Authentifizierung:} Das Ausgangsfrontend enthält Ansätze für
Benutzer- und Authentifizierungsfunktionen. Für die produktive Integration
fehlten jedoch die verbindlichen Schnittstellen und Identitätsdaten der
übergeordneten Plattform. Die erste Iteration arbeitet daher ohne
rollenbasierte Zugriffskontrolle.
\item \textbf{Kubernetes/ArgoCD:} Container-Images und Values-Konfiguration
sind vorbereitet. Ein tatsächliches Deployment war nicht möglich, nachdem
die zuständige Plattformgruppe ihr Projekt nicht fortführte und keine
Zielinfrastruktur zur Integration bereitstellen konnte.
\end{itemize}

\section{Ergebnis, Grenzen und Weiterentwicklung}
\subsection{Erreichter Stand}
Vollständig demonstrierbar sind Aufgaben- und Contentverwaltung, Text- und
Dateiinhalte, Referenzdaten, Darstellungstemplates, geordnete und ungeordnete
Collections, Drag-and-Drop, Varianten, Item-Tags, Validator-Zuordnung sowie
kombinierte Suche und Filter. Datenbank, REST-Backend und Frontend sind über
diese Kernabläufe integriert; der Stand ist automatisiert getestet und lokal
über Docker startbar.

\subsection{Bekannte Grenzen}
\begin{itemize}
\item Validatoren und Modifier werden gespeichert, aber nicht ausgeführt.
\item Für Modifier existiert kein vollständiger Verwaltungseditor.
\item Content-Tags sind backendseitig berücksichtigt, aber nicht als
gleichwertiger eigener UI-Workflow ausgebaut.
\item Die Kompatibilitätsbeziehung \texttt{item\_content\_types} schränkt die
Content-Auswahl noch nicht automatisch ein.
\item Das Schema erlaubt die Wiederverwendung eines Contents in mehreren
Items; die Oberfläche bietet dafür noch keinen eigenen Auswahlablauf.
\item Die Baumansicht behandelt die Struktur in der ersten Iteration als
eindeutige Position, obwohl das relationale Schema Mehrfachzuordnungen
zulässt.
\item Zentrale Authentifizierung, produktives Deployment und externe
Regelausführung fehlen wegen der beschriebenen Plattformabhängigkeiten.
\item Ein Ende-zu-Ende-Test gegen eine echte PostgreSQL-Instanz ist noch offen.
\end{itemize}

\subsection{Priorisierte Verbesserungsmöglichkeiten}
\begin{enumerate}
\item Integrationstests mit PostgreSQL über Testcontainers,
\item zentrale Authentifizierung und Rollen für Autoren,
\item vollständiger Modifier- und Content-Tag-Workflow,
\item automatische Kompatibilitätsfilter für Content-Typen,
\item explizite UI für die Wiederverwendung von Contents und
Mehrfachzuordnungen,
\item Anbindung eines externen Validator-/Modifier-Dienstes,
\item danach eine klar definierte und evaluierbare KI-Unterstützung,
\item produktives ArgoCD-Deployment, sobald Zielplattform und Zugang bestehen.
\end{enumerate}

\section{Fazit}
Das Projekt hat das vorgegebene Schema nicht vereinfacht oder umgangen,
sondern seine zentralen fachlichen Beziehungen in ein nutzbares
Autorenwerkzeug übersetzt. Besonders die Trennung von Item und Content, die
Unterscheidung von Collections und Varianten sowie die konsistente
Positionsverwaltung gehen über einfache CRUD-Masken hinaus.

Nicht erreichte Punkte werden transparent von den implementierten Funktionen
getrennt. Der bestehende Stand bildet eine belastbare erste Iteration:
Kernabläufe sind über Datenbank, Backend und Frontend integriert, getestet und
containerisiert. Die offenen Erweiterungen können auf den vorhandenen
Entitäten und Schnittstellen aufbauen, sobald die fachlichen beziehungsweise
plattformseitigen Voraussetzungen vorliegen.

\begin{thebibliography}{9}
\bibitem{vue} Vue.js Documentation. \url{https://vuejs.org/}
\bibitem{vuetify} Vuetify Documentation. \url{https://vuetifyjs.com/}
\bibitem{pinia} Pinia Documentation. \url{https://pinia.vuejs.org/}
\bibitem{spring} Spring Boot Reference Documentation.
\url{https://docs.spring.io/spring-boot/}
\bibitem{postgres} PostgreSQL 16 Documentation.
\url{https://www.postgresql.org/docs/16/}
\bibitem{docker} Docker Documentation. \url{https://docs.docker.com/}
\bibitem{k8s} Kubernetes Documentation. \url{https://kubernetes.io/docs/}
\bibitem{intern} THM E-Learning-Plattform: bereitgestelltes Frontend,
Datenschema, Lastenheft und interne Abstimmungen.
\end{thebibliography}

\appendix
\section{Anhang}
\subsection{Lokale Ausführung}
\begin{verbatim}
docker compose up --build
# Frontend: http://localhost:8085
# Backend:  http://localhost:8080
# Dummy:    cd frontend && npm run dev:dummy
\end{verbatim}

\subsection{Abschließende Nachweise}
Vor der finalen PDF-Erstellung sind zu ergänzen:
\begin{itemize}
\item Namen, Gruppennummer und konkrete Aufgabenverteilung,
\item die fünf Anwendungsscreenshots unter
\texttt{docs/screenshots/01} bis \texttt{05},
\item Screenshot des GitHub Projects als
\texttt{06-github-project.png},
\item optional ein aussagekräftiger Pull-Request als
\texttt{07-pull-request.png},
\item ein lesbares ER-Diagramm und der endgültige Repository-Link.
\end{itemize}
\end{document}

